% !TeX root = ../main.tex

\chapter{Related Work}\label{chapter:related}

This chapter surveys existing approaches to bufferbloat mitigation, traffic classification, adaptive QoS, and bio-inspired networking, positioning MycoFlow relative to each area.

\section{AQM Evolution}

Proportional Integral controller Enhanced (PIE)~\cite{pan2013pie} achieves CoDel-equivalent performance with lower computational cost by modeling the queue delay as a control system. FQ-PIE combines PIE with per-flow queuing. CAKE~\cite{cake2018} remains the most comprehensive open-source AQM for home gateways, integrating shaping, fair queuing, and DiffServ. None of these systems perform automatic per-device DSCP classification; they all require pre-classified traffic to deliver tier-specific service.

\section{Traffic Classification}

\subsection{Deep Learning Approaches}

GGFAST~\cite{ggfast2023} applies payload-agnostic deep learning to encrypted traffic fingerprinting at line rate; it achieves high accuracy but requires offline model training infrastructure and cloud-side catalog updates, making continuous adaptation incompatible with embedded routers that lack the compute and storage for model inference.

Lotfollahi et al.~\cite{lotfollahi2020deep} use deep learning on raw packet features, achieving high accuracy on encrypted traffic. Again, the resource requirements (GPU inference, gigabytes of training data) preclude embedded deployment on a 256~MB router.

\subsection{Deep Packet Inspection}

nDPI~\cite{ndpi2022} provides open-source DPI classification but requires payload parsing, raising privacy concerns and imposing significant CPU cost on constrained hardware. While an OpenWrt package exists for nDPI, its per-packet payload inspection overhead is incompatible with constrained single-core SoCs running concurrent Wi-Fi, NAT, and AQM tasks.

MycoFlow deliberately avoids both ML and DPI: its heuristic decision tree runs in under 1~$\mu$s per device per cycle on a Cortex-A53 core, requiring no model files and no payload access.

\section{Adaptive QoS and SQM}

SDN-based per-application QoS systems achieve fine-grained traffic differentiation but require a network controller external to the CPE, adding deployment complexity incompatible with consumer home settings where no IT staff are available.

OpenWrt \texttt{qosify} provides static DSCP mapping via nftables rules with fixed port-to-class tables, but no runtime sensing or adaptation. It offers a partial solution for technically skilled users who manually maintain port tables, but cannot respond to changes in traffic patterns or network conditions.

SQM (Smart Queue Management, \texttt{luci-app-sqm}) is the canonical OpenWrt QoS solution~\cite{cake2018}: it configures CAKE or fq\_CoDel with user-specified bandwidth limits and applies AQM. However, SQM has three structural limitations that MycoFlow addresses:

\begin{enumerate}
  \item \textbf{Static bandwidth cap:} SQM requires the operator to manually configure the WAN bandwidth cap once at installation; it does not adjust the cap in response to observed congestion.
  \item \textbf{No automatic classification:} If a user wants gaming traffic to enter CAKE's Voice tin, they must manually assign DSCP marks on each game client---an operation requiring per-device network configuration knowledge.
  \item \textbf{No per-flow distinction:} SQM cannot distinguish between multiple applications running simultaneously on the same host (e.g., a gaming session and a background torrent), because it operates at the qdisc level with no per-flow state.
\end{enumerate}

MycoFlow is designed as a \emph{zero-touch complement} to SQM: it deploys CAKE identically but adds autonomous sensing, classification, and actuation so that no manual DSCP tagging or bandwidth tuning is ever required.

\section{eBPF for Edge Telemetry}

Magnani et al.~\cite{magnani2022} demonstrate adaptive eBPF-based monitoring pipelines on edge routers, showing that kernel-side counters can expose per-flow statistics with negligible overhead. Their work validates the feasibility of eBPF for in-router telemetry and informs MycoFlow's passive TCP RTT observer.

MycoFlow uses eBPF for passive TCP RTT sampling via a TC~clsact hook but primarily relies on \texttt{/proc/net/nf\_conntrack} for flow statistics, as eBPF kernel headers are not available on all OpenWrt builds.

\section{Bio-Inspired Networking}

\subsection{Ant Colony Optimization}

Ant Colony Optimization~\cite{dorigo1996ant} has been applied to routing in multi-hop networks. The core idea---that distributed agents following simple local rules can produce globally optimal paths---maps to MycoFlow's per-device decision making. However, ACO requires agent communication infrastructure absent in home routers; MycoFlow achieves similar local-sensing behavior with purely passive observation.

\subsection{Mycelial Network Research}

Adamatzky and colleagues~\cite{adamatzky2022logics} characterize mycelial networks as fault-tolerant distributed logic gates, demonstrating that biological neural-like computation can emerge from simple local chemical signaling.

Fricker et al.~\cite{fricker2017mycelium} model nutrient transport in mycelial networks as an adaptive flow optimization problem, showing that transport channels in mycelium operate with hysteresis: a sustained signal above a threshold triggers remodeling, preventing oscillation.

Fukasawa et al.~\cite{fukasawa2023foraging} analyze how foraging heuristics in mycelium maximize resource extraction under uncertainty, finding that hysteresis-controlled redistribution outperforms reactive strategies in dynamic environments.

MycoFlow is the first system to apply mycelial foraging and transport principles specifically to the QoS control plane of an embedded router.

\section{Positioning}

Table~\ref{tab:comparison} summarizes the comparison of related approaches against MycoFlow. MycoFlow occupies a unique design point: it runs entirely on the CPE, requires no cloud or SDN controller, uses no ML or DPI, and automatically classifies and steers per-device traffic with runtime adaptation.

\begin{table}[h]
\caption{Comparison of Related QoS Approaches. Emb.=runs fully on CPE; Auto=automatic classification; No~ML=no ML model required; No~DPI=no payload inspection; Adapt=runtime adaptation.}
\label{tab:comparison}
\centering
\small
\begin{tabular}{p{2.2cm}ccccc}
\toprule
\textbf{System} & \textbf{Emb.} & \textbf{Auto} & \textbf{No ML} & \textbf{No DPI} & \textbf{Adapt} \\
\midrule
SQM/CAKE     & \checkmark & $\times$ & \checkmark & \checkmark & $\times$ \\
qosify       & \checkmark & $\times$ & \checkmark & \checkmark & $\times$ \\
nDPI         & \checkmark$^\dagger$ & \checkmark & \checkmark & $\times$ & $\times$ \\
GGFAST       & $\times$   & \checkmark & $\times$   & \checkmark & $\times$ \\
SDN-QoS      & $\times$   & \checkmark & \checkmark & \checkmark & \checkmark \\
\textbf{MycoFlow} & \checkmark & \checkmark & \checkmark & \checkmark & \checkmark \\
\bottomrule
\multicolumn{6}{l}{\footnotesize $^\dagger$nDPI imposes significant CPU cost on constrained hardware.}
\end{tabular}
\end{table}
